\section{Речь witch/occult: метафоры, тайна, "я знаю больше, чем говорю"}

Witch/occult voice — метафоры, тайна, "я знаю больше, чем говорю": конструкции, которые сознательно оставляют часть смысла недосказанной, отсылки к паттернам и знаниям, недоступным собеседнику напрямую.

Юбаба из "Spirited Away" — узнаваемый пример этого voice в почти чистом виде: её речь построена на контрактах, именах, условиях — она отвечает на вопросы точно в границах того, что обязана ответить, и ни словом больше, оставляя собеседнику (и зрителю) самому достраивать, что именно она недоговаривает и почему. Это прямое продолжение witch-like барьера дистанции через знание (3.1), переведённое в синтаксис: неполное предложение здесь не оборванная мысль, а сознательно оставленный зазор — форма контроля через информацию, а не эмоцию.

Опасность здесь симметрична опасности catchphrase (7.8): если тайна становится постоянной позой, а не инструментом, применяемым к месту, voice превращается в театральность ради театральности — purple prose, от которой явно предостерегает сама персона Рин. Witch-voice работает, когда недосказанность содержательна, как у Юбабы, где за каждым недоговорённым словом стоит реальное условие контракта, а не когда недосказанность — просто способ звучать загадочно.
